\documentclass[10pt,a4paper]{amsart}
\usepackage[margin=19mm]{geometry}
\usepackage{amsmath,amssymb,mathtools}
\usepackage[scr=rsfs,cal=euler]{mathalfa}
\usepackage{xfrac}
\usepackage[unicode,hidelinks,bookmarks=false]{hyperref}
\newcommand{\ie}{i.e.\ }
\newcommand{\cf}{cf.\ }
\newcommand{\resp}{respectively\ }
\newcommand{\Subsection}{\subsection}
\providecommand{\nobreakdash}{\nobreak\hbox{-}}
\setlength{\emergencystretch}{2em}
\multlinegap0pt
\usepackage{bbm}
\usepackage{algorithm}\usepackage{algpseudocode}
\usepackage{amssymb}
\usepackage{bm}
\usepackage{caption}
\usepackage[shortlabels]{enumitem}
\usepackage{mathtools}
\newtheorem{definition}{Definition}
\newtheorem{proposition}{Proposition}
\newcommand{\E}{\mathbb E}
\newcommand{\Ind}[1]{\mathbbm{1}{\{#1\}}}
\newcommand{\N}{{\mathbb N}}
\newcommand{\Prob}{\mathbb{P}}
\newcommand{\R}{{\mathbb R}}
\newcommand{\ind}[1]{\mathbbm{1}_{\{#1\}}}
\newcommand{\rd}{\mathrm{d}}
\newcommand{\re}{{\mathrm e}}
\newcommand{\sC}{{\mathscr C}}
\newcommand{\sG}{{\mathscr G}}
\newcommand{\sL}{{\mathscr L}}
\newcommand{\sT}{{\mathscr T}}
\newcommand{\sU}{{\mathscr U}}
\title{The critical percolation window in growing random graphs}
\author{Joost Jorritsma, Pascal Maillard, Peter M\"orters}
\date{Source: arXiv:2512.18937v1 (CC BY 4.0)}
\begin{document}
\maketitle

\noindent\emph{Source and licence.} The critical percolation window in growing random graphs. Version arXiv:2512.18937v1 (CC BY 4.0), distributed by arXiv under the Creative Commons Attribution 4.0 International licence, http://creativecommons.org/licenses/by/4.0/, as recorded in the arXiv licence metadata for this exact version and shown on the abstract page https://arxiv.org/abs/2512.18937v1. Full licence notice: \texttt{licenses/arxiv-2512.18937-CC-BY-4.0.txt}.\par\smallskip

\noindent\textit{Editorial scope.} Counted unit: the article's proposition prop:domination-upper (source0.tex lines 635--645). The article's complete original proof of it is delivered with it (source0.tex lines 646--815, one \texttt{proof} environment of 170 lines). No mathematics is written, repaired or re-derived: the statement and the proof are the article's own lines, in the article's own notation, and no retained line is substituted or re-flowed.  Retained uncounted as the source-local context the proof needs: lines 279--287; lines 498--501; lines 572--587; lines 608--610; lines 627--629.  Nothing was omitted from any retained span, so the package carries no omission record.  No retained line contains a citation, so the wrapper carries no bibliography and no bibliography entry.  No retained line contains a figure environment, an image-inclusion command or drawing code, so no image is delivered and the package declares no asset dependency.  Editorial formatting only, in addition: an AMS \texttt{amsart} preamble that restates the article's own declarations and macros used by the retained lines, namely the environments \texttt{definition}, \texttt{proposition} on separate counters, together with the standard packages those lines need.  The retained environments print in this document as Definition 1, Proposition 1 in order of appearance, whereas the article prints the same items with the numbering of its own full document.  The complete native source of the article as distributed in the arXiv e-print for this exact version is bundled unchanged as \texttt{sources/191435/source0.tex}.
\begin{definition}[$\gamma$-growing random graph]\label{def:pa-ann}
	Let $\gamma\in(-\infty,1)$ and $\beta>0$. Let $\sG_1$ be the graph with a single vertex labeled 1. At discrete steps  $j\in\N\setminus\{1\}$ a new vertex with label $j$ arrives and connects independently by an edge to present vertices $i\in[j-1]:=\{1,\ldots,j-1\}$  with probability
    \begin{equation}\label{eq:conn-prob2}
     p_{ij} = \big(\beta i^{-\gamma}j^{\gamma-1}\big)\wedge 1.
    \end{equation}
	to form $\sG_j$ from $\sG_{j-1}$. 
    We write $\Prob^{\beta}$ and $\E^\beta$ for the law and expectation of the random graph, and omit in our notation the dependency on $\gamma$.
    We write $\sC_n(v)$ for the connected component containing $v$ in $\sG_n$ and $\sL_n(\beta)$ or $\sL_n$ for the largest connected component in $\sG_n$ with arbitrary tie-breaking rule. 
    \end{definition}

\begin{equation}
	\label{eq:mu_beta}
	\mu_\beta(\rd y) := \beta\cdot \big(\re^{(1-\gamma)y}\ind{y<0}+ \re^{\gamma y}\ind{y\ge 0}\big)\,\rd y = \beta \re^{y/2 - (1/2-\gamma)|y|}\,\rd y = \beta \re^{y/2 - 2\beta_c|y|}\,\rd y,
\end{equation}

\begin{equation}\label{eq:i-upper}
\begin{aligned}
x^+_i&=\log(2i-1), 
	&I_i^+&=
    \begin{dcases}
        \big(\log(2i-3), \log(2i-1)\big],\phantom{..}&\text{if }i\ge2,\\
        \{0\},&\text{if }i=1,
    \end{dcases}\\
    \ell^+(x)&=\begin{dcases}\lceil \tfrac{1}{2}\re^x\rceil+1,&\text{if }x>0,\\
    1,&\text{if }x=0,\end{dcases}
    &I^+_{[m,n]}&=\begin{dcases}
        \big(\log(2m-3),\log(2n-1)\big],&\text{if }m\ge2,\\
       \big [0,\log(2n-1)\big],&\text{if }m=1.
    \end{dcases}
    \end{aligned}
\end{equation}

\begin{definition}[$\gamma$-branching random walk]\label{def:pa-brw}
	Fix $\gamma\in[0, 1/2)$, and let $\beta>0$. Let $\xi$ be the random sequence of atoms in increasing order of a Poisson point process on $\R$ with intensity $\mu_\beta$ defined in~\eqref{eq:mu_beta}. The $\gamma$-branching random walk $(X_s)_{s\in\mathcal U}$ starting from $x\in\R$ is defined to be the branching random walk with reproduction law $\mathcal L_\xi$ starting from $x$. We denote by $\Prob_x^\beta$ its law and write $\E_x^\beta$ for the expectation with respect to it.
\end{definition}

\begin{definition}[Real, fake, and colliding particles]\label{def:real}
Let \smash{$\sT^\pm_{[m,n]}(v)$} be the total progeny of the $\gamma$-branching random walk restricted to \smash{$I^\pm_{[m,n]}$}.  We declare the root to be \emph{real}, and assign types to all other particles in \smash{$\sT^\pm_{[m,n]}$} recursively in the order $\prec$.  If a particle $t$ has position \smash{$X_t\in I^\pm_i$} for a vertex $i\in[m,n]$, then $t$ is called \emph{real} if its parent $t'$ is real and if there is no real particle \smash{$s\in \sT^\pm_{[m,n]}$} with $s\prec t$ and \smash{$X_s\in I^\pm_i$}.  If \smash{$I^\pm_i$} contains a real particle $s\prec t$ and $t'$ is real, we call $t$ both \emph{colliding} and \emph{fake}.  All descendants of these colliding particles  are also called fake.  We write \smash{$\mathrm{Real}(\sT^\pm_{[m,n]}(v))$} and \smash{$\mathrm{Fake}(\sT^\pm_{[m,n]}(v))$} for the sets of real and fake particles, respectively, and we denote the set of colliding particles with position in \smash{$I^\pm_i$} by \smash{$\mathrm{Colliding}(\sT^\pm_{[m,n]}(v),I_i^\pm)$}. Note that \smash{$\mathrm{Colliding}(\sT^\pm_{[m,n]}(v),I^\pm_i)\subseteq \mathrm{Fake}(\sT^\pm_{[m,n]}(v))$} for every $i$. 
\end{definition}

\begin{proposition}[Coupling the branching random walk with a component]\label{prop:domination-upper}
Let $\beta\in(0,1/2]$ and $\gamma\in[0,1)$. Consider the  $\gamma$-growing random graph from Definition~\ref{def:pa-ann}, and the $\gamma$-branching random walk from Definition~\ref{def:pa-brw}. Assume first that $m, v, n\in\N$ such that $m\le v\le n$. There exists a coupling of $\sT^-_{[m,n]}(v)$ and $\sC_{[m,n]}(v)$ such that on this coupling, for all $k\in\N$,
\begin{equation}\label{eq:coupling-lower}
\big\{\ell^-(X_s): s\in\mathrm{Real}\big(\sT^-_{[m,n]}(v)\big), |s|\le k\big\}\, \subseteq\, 
			\{j\in \sC_{[m,n]}(v): d_{[m,n]}(v, j)\le k\}.
\end{equation}
Assume next that $v\le n$, and $m\in[1+\mathbbm1_{v\ge2}, v]$.
There exists a coupling of $\sT^+_{[m,n]}(v)$ and $\sC_{[m,n]}(v)$ such that on this coupling, for all $k\in\N_0$,	
	\begin{equation}\label{eq:coupling-upper}
			\{j\in \sC_{[m,n]}(v): d_{[m,n]}(v, j)\le k\}                                       \, \subseteq\,  \big\{\ell^+( X_s): s\in\mathrm{Real}\big(\sT^+_{[m,n]}(v)\big), |s|\le k\big\}.\end{equation}
\end{proposition}

\begin{proof}
	Fix $m\le v\le n$. We construct a probability space on which we simultaneously sample the killed branching random walk \smash{$\sT_{[m,n]}^\pm(v)$} and the component $\sC_{[m,n]}(v)$. We use the construction both for the lower and upper bound. To do so, we use the following collections of iid uniform random variables: 
    \begin{equation}
        U_{s, j}\sim\mathrm{Unif}[0,1], \quad s\in \sU,\  j\in[m,n], \qquad \text{and}\quad R_{i,j}\sim \mathrm{Unif}[0,1], \quad i,j\in[m,n], \ i<j.
    \end{equation}
     We construct a sample of the branching random walk killed outside $\smash{I_{[m,n]}^\pm}$ iteratively. We start the branching random walk with the root positioned at $\log v$ or $\log(2v-1)$ for lower and upper bound, respectively.
    Let $q_{s,j}^\pm$ be the probability that $s\in\sU$, given its position $X_s$, produces at least one offspring in the interval \smash{$I_{j}^\pm$}. Then we may couple the branching random walk with the uniform random variables such that for all $s\in \sU$ and $j\in[m,n]$,
    \[
    \Ind{\exists i\in\N: X_{si}\in I_j^\pm}=\Ind{U_{s,j}\le q_{s, j}^\pm}.
    \]
    Given the sample of the branching random walk and the uniform random variables, we determine the component of $\sC_{[m,n]}(v)$ via the order $\prec$ of branching random walk particles. We call particles in the branching random walk real and fake following Definition \ref{def:real}. Each interval $\smash{I_i^\pm}$ contains at most one real particle. Assume first that $\smash{I_i^\pm}$ contains a real particle, denoted by $s_i$. 
    Then we couple the edges in $\sG_{[m,n]}$ by setting 
    \begin{equation}\label{eq:coupling-edges}
    \Ind{i\sim j}=\Ind{U_{s_i, j}\le p_{ij}}, \qquad
      \parbox[t]{0.6\linewidth}{if there exists a real particle $s_i$ such that  $X_{s_i}\in I_i^\pm$, \\and there is no real particle $s_j$ such that $X_{s_j}\in I_j^\pm$ with $s_j\prec s_i$.}
    \end{equation}
    Edges $\{i,j\}$ in $\sG_{[m,n]}$ are now uniquely determined provided that $I_i^\pm$ or $I_j^\pm$ contains a real particle in \smash{$\sT^\pm_{[m,n]}$}, as the offspring of the smallest real particle in the ordering decides presence of the edge.
    The only undetermined edges $\{i,j\}$ in $\sG_{[m,n]}$ are those such that  $I_i^\pm\cup I_j^\pm$ contains no real particle. We set, for $i<j$,
    \[
    \Ind{i\sim j}=\Ind{R_{i,j}\le p_{ij}}, \qquad \text{if both $I_i^\pm$ and $I_j^\pm$ contain no real particles}.
    \]
    Note that the described coupling indeed preserves the marginals for the connected component and killed branching random walk. We now prove the lower bound \eqref{eq:coupling-lower}.

    \medskip\noindent\emph{Lower bound.\ }
    On the described coupling, provided that $q_{t,j}^-\le p_{ij}$ for each real particle $t\in I_i$, for all $i, j\in[m,n]$, we have the following implication:  If a particle $s$ in generation $k$ is real, then its ancestral line towards the root has a corresponding path in $\sC_{[m,n]}$ by projecting particle positions on the ancestral path to vertices using $\ell^-(x)=\lfloor\re^x\rfloor$. We show next that $\smash{q_{t,j}^-}\le p_{ij}$ is satisfied for all particles $t$. We recall that $\smash{q_{t,j}^-}$ denotes the probability that a particle $t$ produces at least one offspring in $\smash{I_{j}^-}$. Bounding $p_{ij}\ge1-\exp(-p_{ij})$, it suffices to show that, for $t$ with $X_t\in I_i^-$,
	\begin{equation}\label{lower}
		q_{t,j}^-= 1-\exp\big(-\mu_\beta
		\big( (\log j -X_{t}, \log (j+1) -X_{t}]\big)\big) \le 1-\exp\big(\beta(i\vee j)^{\gamma-1}(i\wedge j)^{-\gamma}\big).
	\end{equation}
	By definition of 
    $\mu_\beta$ in~\eqref{eq:mu_beta}, $q_{t,j}^-$ is maximized when $X_{t}$ is at the leftmost value in $I^-_i=[\log i, \log(i+1))$. Thus, 
	\begin{align*}
		\mu_\beta
		\big( (\log j -X_{t}, & \log (j+1) -X_{t} ]\big)
		\leq \mu_\beta \big( (\log (j/i), \log ((j+1)/i) ]\big).
	\end{align*}
	If $i<j$ and $\gamma\neq 0$, then
	\begin{align*}
		\mu_\beta \big( (\log (j/i), \log ((j+1)/i) ]\big)
		 & = \beta \int_{\log (j/i)}^{\log ((j+1)/i)}
		\re^{\gamma y} \, \rd y
		= \tfrac\beta\gamma \Big(((j+1)/i)^\gamma-(j/i)^\gamma \Big)                                                 \\
		 & =\tfrac\beta\gamma  \big((1+1/j)^\gamma-1 \big) j^\gamma i^{-\gamma} \leq \beta j^{\gamma-1} i^{-\gamma}.
	\end{align*}
    The last bound is easily verified by substituting $x=1/j$, then verifying the bound at $x=0$, and observing that the derivative in $x$ of the left-hand side is at most the derivative of the right-hand side. When $i<j$ and $\gamma=0$, 
    \[
    \mu_\beta \big( (\log (j/i), \log ((j+1)/i) ]\big) = \beta\big(\log((j+1)/i)-\log(j/i)\big)=\beta\log(1+1/j)\le \beta/j.
    \]
    	Similarly, if $i>j$, then for any $\gamma<1$
	\begin{align*}
		\mu_\beta \big( (\log (j/i), \log ((j+1)/i) ]\big)
		 & = \beta \int_{\log (j/i)}^{\log ((j+1)/i)}
		\re^{(1-\gamma) y} \, \rd y
		= \tfrac\beta{1-\gamma} \Big(((j+1)/i)^{1-\gamma}-(j/i)^{1-\gamma} \Big)            \\
		 & =\tfrac\beta{1-\gamma}  \big((1+1/j)^{1-\gamma}-1 \big) j^{1-\gamma}i^{\gamma-1}
		\leq \beta i^{\gamma-1} j^{-\gamma} .
	\end{align*}
    These three cases prove~\eqref{lower} and the lower bound \eqref{eq:coupling-lower} follows.

\pagebreak[3]

\medskip\noindent\emph{Upper bound.\ } We next prove \eqref{eq:coupling-upper} via the same coupling. Recall that in the upper bound embedding the interval $I_1^+$ is the singleton $\{0\}$. Since we consider the component of $v$ in the induced subgraph on $[m,n]$ with $m\ge 1+\mathbbm{1}_{v\ge 2}$, vertex~$1$ does not appear in $\sC_{[m,n]}(v)$ whenever $v\ge 2$. In the remaining case $v=1$, the root is a real particle in $\smash{I_1^+}=\{0\}$, and $1$ is contained in the set on the right-hand side in~\eqref{eq:coupling-upper}. Thus, it suffices to show that 
\begin{equation}\label{eq:coupling-upper2}
\{j\in \sC_{[m,n]}(v): d_{[m,n]}(v, j)\le k, j\ge 2\}                                       \, \subseteq\,  \big\{\ell^+( X_s): s\in\mathrm{Real}\big(\sT^+_{[m,n]}(v)\big), |s|\le k\big\}.
\end{equation}
We show below that $p_{ij}\le \smash{q^+_{t,j}}$ for all real particles $t\in \smash{I^+_i}$ and all $i, j\in[m,n]$ and $j\ge 2$. Under this assumption, we claim that \eqref{eq:coupling-upper2} follows by induction on~$k$. For $k=0$, it is trivial, as both sets consist only of the vertex $v$. Assuming the inclusion holds for $k$, we advance the induction and argue next that it also holds for $k+1$. 
We decompose
\[
\begin{aligned}
\big\{j\in\sC_{[m,n]}(v) &: d_{[m,n]}(v,j)\le k+1\big\}\\ 
&= \big\{j\in\sC_{[m,n]}(v) : d_{[m,n]}(v,j)\le k\big\} \cup \big\{j\in\sC_{[m,n]}(v) : d_{[m,n]}(v,j)= k+1\big\}.
\end{aligned}
\]
By the induction hypothesis, the first set on the right-hand side is a subset of the real particles in the first $k$ generations. So, we have to argue that 
\begin{equation}\label{eq:upper-coupling-inclusion} \big\{\ell^+( X_s): s\in\mathrm{Real}\big(\sT^+_{[m,n]}(v)\big), |s|\le k+1\big\}\supseteq \big\{j\in\sC_{[m,n]}(v) : d_{[m,n]}(v,j)= k+1, j\ge 2\big\}.
\end{equation}
Thus, we consider the event that a breadth-first exploration of $\sC_{[m,n]}(v)$ discovers a vertex $j\in[m,n]$ exactly at generation $k+1$. In particular, we work on the event that $d_{[m,n]}(v,j)>k$. For proving \eqref{eq:upper-coupling-inclusion}, we may assume without loss of generality that there is no real particle $s$ in the first~$k$ generations such that $\ell^+(X_s)=j$. Let 
$$\mathrm{Real}_k(v):=\big\{\ell^+(X_s) : |s|=k, s\in\mathrm{Real}\big(\smash{\sT^+_{[m,n]}}(v)\big)\big\}.$$
By \eqref{eq:coupling-edges}, presence of an edge $\{i,j\}$ is decided  by the offspring of the smallest real particle in $\smash{I_i^+ \cup I_j^+}$ with respect to the ordering $\prec$. There are no real particles in the union of $\big(\smash{I_i^+} : i\in\mathrm{Real}_k(v)\big)$ before generation $k$ by definition of real particles, and $\smash{I_j^+}$ contains no real particle before generation $k+1$ by assumption. Therefore, none of the edges $\{i, j\}$ with $i\in\mathrm{Real}_k$ has been determined before generation $k+1$. These edges are revealed exactly when we expose the offspring of the real particles in generation~$k$.

\smallskip 
If none of the real particles in generation $k$ produces an offspring in $\smash{I^+_j}$, then $U_{s,j}> \smash{q^+_{s,j}}$ for all real particles~$s$ in generation $k$. As we assume that $\smash{q^+_{s,j}}\ge p_{ij}$ for all real particles~$s$ with $X_s\in \smash{I^+_i}$ and all $i, j\in[m,n]$ with $j\ge 2$, all edges between $\mathrm{Real}_k$ and $j$ are absent on this event, and we necessarily have that $d_{[m,n]}(v, j)\ge k+2$ as we work on the event that $d_{[m,n]}(v,j)>k$.

\smallskip 
Therefore, the only way that $d_{[m,n]}(v,j)$ can be discovered by the breadth-first exploration in generation $k+1$ while none of the real particles up to generation $k$ project to $j$, is if there is a smallest real particle $s$ in generation $k$ (with respect to $\prec$) that produces an offspring $si$ with position $X_{si}\in \smash{I^+_{j}}$. The first offspring that $s$ produces in $\smash{I^+_j}$ must be real because $\smash{I^+_j}$ did not contain a real particle before.
The containment in  \eqref{eq:upper-coupling-inclusion} follows provided that $\smash{q^+_{t,j}}\ge p_{ij}$ for all real particles $t\in \smash{I^+_i}$ and all $i, j\in[m,n]$ with $j\ge 2$. Thus,~\eqref{eq:coupling-upper2} and~\eqref{eq:coupling-upper} follow under the same condition.  

\medskip
\noindent\emph{Verification of the one-step inequality $q^+_{t,j}\ge p_{ij}$.}
We finish the proof of \eqref{eq:coupling-upper} by showing the inequality $\smash{q^+_{t,j}}\ge p_{ij}$ for all real particles $t\in \smash{I_i^+}$ and all $i,j\in[m,n]$ with $j\ge 2$.  Since $\smash{q^+_{t,j}}$ is the probability that particle $t$ produces at least one offspring in $\smash{I_j^+}$, we will show that for any $t$ with $\smash{\ell^+}(X_t)=i$, and $j\ge 2$,
	\begin{equation}\label{upper}
		q^+_{t,j} = 1-\exp\Big(-\mu_\beta
		\big( (\log (2j-3) -X_{t}, \log(2j-1) -X_{t}]\big)\Big) \ge \beta(i\vee j)^{\gamma-1}(i\wedge j)^{-\gamma}.
	\end{equation}
    Note that $j\ge 2$ implies that $\log(2j-3)$ is well-defined.
	By the definition of the measure $\mu_\beta$ in~\eqref{eq:mu_beta}, the left-hand side is minimized when $X_{t}$ is at the rightmost value in $I^+_i$, which is $\log (2i-1)$ for all $i\ge 1$ by~\eqref{eq:i-upper}. Thus,
    \begin{align*}
		\mu_\beta
		\big( (\log (2j-3) -X_{t}, & \log(2j-1) -X_{t} ]\big)
		\geq \mu_\beta \Big( (\log \big(\tfrac{2j-3}{2i-1}\big), \log \big(\tfrac{2j-1}{2i-1}\big) \big]\Big).
	\end{align*}
    We distinguish four cases. 

    \smallskip 
    \noindent\emph{Case 1: $1\le i<j$, $\gamma>0$.\ }
	If $1\le i<j$ and $\gamma\neq 0$, then 
	\begin{align}
		\mu_\beta \Big( (\log \big(\tfrac{2j-3}{2i-1}\big), \log \big(\tfrac{2j-1}{2i-1}\big) \big]\Big)
		 & = \beta \int_{\log \frac{2j-3}{2i-1}}^{\log \frac{2j-1}{2i-1}}
		\re^{\gamma y} \, \rd y
		= \tfrac\beta\gamma \Big(\big(\tfrac{2j-1}{2i-1}\big)^\gamma-\big(\tfrac{2j-3}{2i-1}\big)^\gamma \Big)                                                \label{upper-int} \\
		 & =\tfrac\beta\gamma  (2i-1)^{-\gamma}\Big((2j-1)^\gamma-(2j-3)^{\gamma}\Big) \nonumber\\&= \beta(i-\tfrac12 )^{-\gamma}\cdot\tfrac{1}{\gamma}\Big((j-\tfrac12 )^\gamma-(j-\tfrac32)^{\gamma}\Big).\nonumber
	\end{align}
    We use that $(x^a-(x-1)^a)/a\ge x^{a-1}$ for $x>1$ and $a\in(0,1)$, and that $1-\exp(-x)\ge x-x^2/2$ for $x>0$. This yields,
    \begin{align}
    1-\exp\Big(-\mu_\beta
		&\big( (\log (2j-3) -X_{s}, \log(2j-1) -X_{s}]\big)\Big)\nonumber\\
        &\ge
        1-\exp\Big(\beta(i-\tfrac12 )^{-\gamma}(j-\tfrac12 )^{\gamma-1}\Big) \label{eq:upper-12}\\
        &\ge \beta(i-\tfrac12 )^{-\gamma}(j-\tfrac12 )^{\gamma-1} - \tfrac{\beta^2}{2}(i-\tfrac12 )^{-2\gamma}(j-\tfrac12 )^{2(\gamma-1)}.\nonumber
        \end{align}
    Recall from \eqref{upper} that we aim to show that the right-hand side is at least $\beta i^{-\gamma}j^{\gamma-1}$. Rearranging terms, we obtain that \eqref{upper} is implied if 
    \begin{equation}\nonumber
    \begin{aligned}
    \tfrac{\beta}{2}&\le (i-\tfrac12 )^{\gamma}(j-\tfrac12 )^{1-\gamma}-i^{-\gamma}j^{\gamma-1}(i-\tfrac12 )^{2\gamma}(j-\tfrac12 )^{2(1-\gamma)} \\
    &= \big(\tfrac{i-\frac12 }{i}\big)^{\gamma}\big(\tfrac{j-\frac12 }{j}\big)^{1-\gamma}\Big(i^\gamma j^{1-\gamma} - (i-\tfrac12 )^{\gamma}(j-\tfrac12 )^{1-\gamma}\Big).
    \end{aligned}
    \end{equation}
    The first two factors on the right-hand side are increasing in $i$ and $j$ as $\gamma\in[0,1)$. Their product is at least $\tfrac12 $ as $j>i\ge 1$. Since we assume $\beta\le \tfrac12 $, it suffices to show that 
    \begin{equation}\label{eq:1ijgamma}
    \tfrac12 \le i^\gamma j^{1-\gamma}- (i-\tfrac12 )^{\gamma} (j-\tfrac12 )^{1-\gamma}.
    \end{equation}
    To this end, consider the function 
    $$
    f(x,y) = x^\gamma y^{1-\gamma} - (x-\tfrac12 )^\gamma(y-\tfrac12 )^{1-\gamma}, 
    \text{ for $1\le x\le y$.}  $$
    Its derivative with respect to $y$ is
    \[
    \frac{\partial}{\partial y}f(x,y) = (1-\gamma)\left(\big(\tfrac x y\big)^\gamma - \big(\tfrac{x-\frac12 }{y-\frac12 }\big)^{\gamma}\right) \ge 0,   
    \]
    since $x/y \ge (x-\tfrac12 )/(y-\tfrac12 )$ for $y\ge x\ge 1$. Moreover, we have that
    $
    f(x,x) = x - (x-\tfrac12 ) = \tfrac12 
    $
    for every $x\ge 1$.
    It follows that $f(x,y) \ge \tfrac12 $ for all $y\ge x\ge 1$, and therefore $f(i,j) \ge \tfrac{1}{2}$ for all $1\le i<j$. This proves~\eqref{eq:1ijgamma} when $1\le i<j$, $\gamma>0$.

    \smallskip 
    \noindent\emph{Case 2: $i\!>\!j\ge\! 2$, $\gamma>0$}. Verbatim reasoning applies as in Case 1, with the exponent  $\gamma$ replaced by $1\!-\!\gamma$.

    
    \smallskip 
    \noindent\emph{Case 3: $j>i\ge 1$, $\gamma=0$}. Recalling \eqref{upper-int}, we have in this case, using $\log(1+1/x)\ge 1/(1+x)$ for $x>0$,
    \begin{align*}
    1-\exp\Big(-\mu_\beta \Big( (\log \big(\tfrac{2j-3}{2i-1}\big), \log \big(\tfrac{2j-1}{2i-1}\big) \big]\Big)\Big)
    &= 1-\exp\Big(-\beta\log\big(\tfrac{2j-1}{2j-3}\big)\Big) \\
    &= 1-\exp\Big(-\beta\log\big(1+\tfrac{1}{j-\frac32}\big)\Big) 
    \ge 1-\exp\Big(-\tfrac{\beta}{j-\frac12 }\Big).
    \end{align*}
    Following the same reasoning as in Case 1 from \eqref{eq:upper-12}, the right-hand side is at least $\beta/j$.
    
    \smallskip 
    \noindent\emph{Case 4: $i>j\ge 2$, $\gamma=0$.}
    We use a similar strategy. We compute the integral as in \eqref{upper-int}, with the exponent $\gamma$ replaced by $1-\gamma=1$. We obtain
    \begin{align*}
    1-\exp\Big(-\mu_\beta \Big( (\log \big(\tfrac{2j-3}{2i-1}\big), \log \big(\tfrac{2j-1}{2i-1}\big) \big]\Big)\Big)
    &= 1-\exp\Big(-\beta \big(\tfrac{2j-1}{2i-1}-\tfrac{2j-3}{2i-1}\big)\Big) = 1-\exp\big(-\tfrac{\beta}{i-\frac12 }\big).
    \end{align*}
    Analogously to Case (1) from \eqref{eq:upper-12}, it follows that the right-hand side is at least $\beta/i$ for $i\ge 2$ and $\beta\in(0, \tfrac12 ]$.    This proves~\eqref{upper} for all values $i,j\in[m,n]$ with $j\ge 2$, and finishes the proof of \eqref{eq:coupling-upper}. 
\end{proof}

\end{document}
